% =================
% INTRODUCTION
% -=====--=-=-=-==-=
\section{Introduction}
A growing body of work has identified \textit{evaluation awareness} in large language models (LLMs). That is, LLMs being "aware"--either verbally in chain-of-thought or internally--to themselves being evaluated, thus potentially altering behavior. The most prominent method for identifying internal evaluation awareness involves \textit{linear probes}: training linear classifiers on frozen activations \citep{nguyen2025probing, chaudhary2025evaluation}.

However, given that evaluation and deployment prompts typically have distinct structures, claims of identifying evaluation awareness via linear probes could instead just be identifying a prompt's superficial \textit{format}. In particular, \cite{devbunova2026format} has shown that linear "evaluation awareness" probes seem to mainly be format detectors: in a 2x2 design crossing MCQ format with prompt style, a probe trained to separate benchmark from deployment prompts misclassified 94\% of MCQ-formatted deployment prompts as evaluation.

Because Devbunova's dataset consists of repurposed existing data (MMLU, ARC, GSM8K, ShareGPT, and MuTual), the evaluation/deployment label confounds against \textit{topic} (what a prompt is about). More fundamentally, no prompt ever appears in more than one format (format varies only across different prompts and never within one), so the design can detect the format confound but cannot quantify the format-independent evaluation signal.

Therefore, we extend this study in the following methods:
%  WE'LL DO THIS AT THE END


% =================
% METHODS
% -=====--=-=-=-==-=
\section{Methods}

\subsection{The crossed factorial object}

We adopt the crossed context-by-format design of \citet{devbunova2026format},
which varies usage context independently of surface format, and extend it
from per-cell probe accuracy to an additive decomposition of the hidden
state itself.

We use a crossed $2 \times 2$ design with
stated response purpose, 
$p \in \{\mathrm{evaluation},\mathrm{deployment}\}$
and response format,
$f \in \{\mathrm{benchmark},\mathrm{casual}\}$.
At each layer, let $h_{p, f}$, denote the hidden state, assigning $p=+1$ for evaluation and $p=-1$ for deployment, and $f=+1$ for benchmark and $f=-1$ for casual. We then decompose $h_{p,f}$:

\begin{equation}
h_{p,f} =  pA + fB + pf\,C  ,
\label{eq:decomp}
\end{equation}

where  $A$ is the format-shared purpose effect, $B$ is the format effect, and $C$ is the
purpose-by-format interaction. 

The matched within-format purpose contrasts are
\begin{equation}
\delta_{\mathrm{bench}}
    = h_{EB}-h_{DB}=2(A+C),
\qquad
\delta_{\mathrm{casual}}
    = h_{EC}-h_{DC}=2(A-C).
\end{equation}
The factorial components therefore satisfy
\begin{equation}
A=\frac{\delta_{\mathrm{bench}}+\delta_{\mathrm{casual}}}{4},
\qquad
C=\frac{\delta_{\mathrm{bench}}-\delta_{\mathrm{casual}}}{4},
\qquad
B=\frac{(h_{EB}+h_{DB})-(h_{EC}+h_{DC})}{4}.
\end{equation}
Thus, $A$ is the average difference between evaluation and deployment
across formats, $B$ is the average difference between benchmark and
casual formats across purposes, and $C$ captures how the difference
between evaluation and deployment varies by format. Because all three
quantities are computed within matched payload groups, properties shared
by the four renderings, such as task content and topic, cancel.

\subsection{Dataset and activations}

We constructed 320 purpose-neutral payload blocks covering 16 topics.
For each topic, ten payloads were generated by Grok-4.3 and ten by
Gemini-3.1-Pro. A separate GPT-5.4 validator checked each payload for
answer uniqueness, plausible but incorrect distractors, and equivalence
across response formats. Each block contains a context, a question, one
correct answer, and three distractors. Each payload was rendered using six benchmark-style templates, five
casual templates, and one additional template. The additional template
alternated between the original benchmark-style and casual templates
across payloads. For each payload-template combination, we produced three renderings:
one with an evaluation-purpose statement, one with a
deployment-purpose statement, and one without an explicit purpose
statement.

The resulting dataset contains 11{,}520 items. The primary analysis
uses the 7{,}680 stated-purpose items, corresponding to 3{,}840 matched
evaluation/deployment pairs. The remaining 3{,}840 no-cue items are
analyzed separately and are never pooled with the primary set. Because
the same payload appears across purposes and formats, task content,
topic, and generator are held fixed by construction.

We extract activations from pinned revisions of
Llama-3.1-8B-Instruct, Llama-3.1-70B-Instruct, and
Llama-3.3-70B-Instruct using the same items and item order. Each prompt
is passed through the checkpoint's chat template as a single user
message with a generation prompt appended. At every transformer block,
we record the representation of the final real prompt token after the
block output and before the model's final normalization. Inference uses
bfloat16, and activations are stored in float32. Exact model revisions,
item hashes, and activation-file hashes are reported in the
reproducibility inventory.

\subsection{Probes, splits, and statistical analysis}

At each layer, we fit an $L_2$-regularized logistic-regression probe
with $C=0.1$ on standardized activations. Both the standardizer and the
probe are fitted only on the training data. In each of five folds, we
hold out one purpose-vocabulary family [DEFINED IN APPENDIX], so the
probe is tested on purpose wording not observed during training. We
combine predictions across the five test folds before calculating ROC
AUC and balanced accuracy.

We evaluate four probe regimes. The \textbf{confounded} probe trains on
evaluation prompts in benchmark format and deployment prompts in
casual format, then evaluates all four purpose-by-format cells. The
\textbf{decorrelated} probe trains on all four cells. The
\textbf{benchmark-to-casual} regime trains only on benchmark prompts and tests
on casual prompts, while the \textbf{casual-to-benchmark} regime uses the
reverse direction. [CITE DEVBUNOVA]

As a check for textual shortcuts, we remove the explicit purpose
statement and test whether simple text classifiers can still
distinguish evaluation from deployment prompts. We also compare
mixed-format and single-format probes using equal amounts of training
data at Llama-3.1-8B layer 12. Confidence intervals are calculated by
bootstrapping complete payload blocks, preserving the matched
purpose-by-format structure. Further implementation and control
details are provided in the appendix.










% =================
% RESULTS
% =================
\section{Results}


\subsection{Purpose is strongly decodable but only partially format-invariant}
\label{sec:results_phenomenon}

Across all three models, stated response purpose becomes strongly linearly decodable within the first fifth of the network. 
The decorrelated probe reaches peak ROC AUCs of 
\(0.989\), \(0.998\), and \(0.996\) 
for Llama-3.1-8B, Llama-3.1-70B, and Llama-3.3-70B, respectively, 
before declining gradually toward the final layer (Figure~\ref{fig:depth_trajectory}). 
Purpose information therefore emerges substantially later than surface format, 
but is represented with near-ceiling discriminability at intermediate depth.

High purpose decodability does not imply complete invariance to response format.
We quantify this using the cross-format transfer gap,
\[
g_\ell
=
\mathrm{AUC}_{\mathrm{both},\ell}
-
\frac{
    \mathrm{AUC}_{B\rightarrow C,\ell}
    +
    \mathrm{AUC}_{C\rightarrow B,\ell}
}{2},
\]
where \(B\rightarrow C\) and \(C\rightarrow B\) denote the benchmark-to-casual and casual-to-benchmark probes, respectively. The gap becomes small at
intermediate depth, reaching minima of \(0.0360\), \(0.0104\), and \(0.0123\)
AUC in the three models, respectively. However, this convergence is
non-monotonic: the gap subsequently grows to \(0.0946\), \(0.0578\), and
\(0.0807\) at the final layer.

The resulting late-layer widening is \(0.0585\) AUC for Llama-3.1-8B
(\(95\%\) CI \([0.0523, 0.0647]\)), \(0.0474\) for Llama-3.1-70B
(\([0.0431, 0.0521]\)), and \(0.0684\) for Llama-3.3-70B
(\([0.0630, 0.0732]\)). Thus, purpose becomes highly portable across formats
at intermediate depth, but this portability partially breaks down again near
the output. The dashed erasure trajectories in
Figure~\ref{fig:depth_trajectory} are analyzed in the following subsection.

% ============================================================
% Figure 1: Phenomenon
% ============================================================
\begin{figure*}[t]
    \centering
    \includegraphics[width=\textwidth]{figures/full_layer_depth_trajectory.pdf}
    \caption{
        \textbf{Purpose is strongly decodable but only partially format-invariant.}
        Purpose decodability and the cross-format transfer gap are shown across
        normalized network depth for all three models. Solid lines denote the
        unmodified representations; dashed lines show rank-64 matched-format
        subspace erasure, analyzed in the following subsection. Purpose AUC rises
        rapidly to near-ceiling values at intermediate depth, while the
        cross-format transfer gap contracts and subsequently widens again toward
        the final layers.
    }
    \label{fig:depth_trajectory}
\end{figure*}


\subsection{A low-rank subspace derived from the matched format contrast accounts for a majority of the portability penalty.}



A relatively low-rank format-associated component contributes disproportionately to the failure of linear purpose readouts to transfer across formats, even though most linearly decodable format information remains.

after rank-64 projection, format AUC is still roughly .922, .954, and .958.


within the region where raw joint-format purpose AUC is ≥.95, projection changes joint-format purpose AUC by only about 0.0006–0.0013, while improving both transfer directions by ~.02–.04 AUC.



% ============================================================
% Figure 2: Structure
% ============================================================
\begin{figure*}[t]
    \centering
    \includegraphics[width=\textwidth]{figures/full_layer_gap_erasure.pdf}
    \caption{
        \textbf{The cross-format transfer gap is concentrated in a
        matched-format subspace.}
        Cross-format transfer gap across normalized depth under increasing ranks
        of matched-format PCA erasure, together with the network-wide integrated
        gap and fraction removed. Increasing erasure rank systematically reduces
        the portability gap across all three models.
    }
    \label{fig:gap_erasure}
\end{figure*}



\subsection{A low-rank format-associated subspace explains most of the cross-format transfer gap}
\label{sec:results_projection}

We quantify the cross-format transfer gap's association with format-related variation. 
For each held-out purpose fold, we construct a basis from the matched format %ian will hate me% 
vectors by removing the leading rank-n subspace from activations, refitting the purpose probes on the new projected representations. 
This is done only on the training folds, creating a subspace tKTKTKTKTKT

Accorss all three models,  we remove this format-associated subspace, reducing the cross-format transfer gap as rank increases (\ref{fig:gap_erasure}).
Integrating this gap over normalized network depth, the 64 rank projection reduces this gap from \(0.0793\) to \(0.0329\) for Llama-3.1-8B, from \(0.0550\) to \(0.0205\) for Llama-3.1-70B, and from \(0.0631\) to \(0.0219\) for Llama-3.3-70B. These correspond to reductions of \(58.5\%\), \(62.8\%\), and \(65.3\%\), respectively. The reduction is monotonic with projection rank across checkpoints \ref{fig:gap_erasure}.

{\color{blue}
Crucially, rank-64 projection substantially improves cross-format transfer while
leaving joint-format purpose decodability essentially unchanged. The projection
also does not eliminate format information itself: after rank-64 projection,
format remains strongly linearly decodable, with median ROC AUCs of \(0.922\),
\(0.954\), and \(0.958\) for Llama-3.1-8B, Llama-3.1-70B, and
Llama-3.3-70B, respectively. Thus, the reduction in the transfer gap is not
explained by simply removing all linearly decodable format information.

The effect is substantial but incomplete. Rank-64 projection leaves
approximately \(35\%\)--\(42\%\) of the integrated cross-format transfer gap
intact. We therefore conclude that a low-rank format-associated component
accounts for a majority of the measured transfer gap, but not all of it.


\subsection{Matched-format interchange controls downstream internal format state}
\label{sec:results_interchange}

The projection results identify a low-rank subspace associated with the
cross-format transfer gap, but do not establish whether that subspace is
functionally used during the forward computation. We therefore perform
interchange patching on the rank-16 matched-format subspace. For each
same-purpose matched pair, we replace the source representation's projection
onto this subspace with the corresponding projection from the opposite-format
counterpart, then continue the forward pass normally. We compare this structured
intervention with a random perturbation matched in magnitude to the actual swap.

Across all three checkpoints, the structured interchange reliably drives the
downstream internal format readout toward the opposite-format counterpart
(Figure~\ref{fig:interchange}). Interchange-intervention accuracy rises from near
chance in early layers to approximately \(0.99\) at later depth for both 70B
models, with a similar trajectory in Llama-3.1-8B. In contrast, the
magnitude-matched random control remains substantially lower. Thus, the
matched-format-associated subspace is not merely correlated with format at the
injection layer: intervening on it systematically changes the format information
available downstream.

The effect on the purpose readout is strongly depth-dependent. Structured
interchange produces greater purpose-readout disruption than the matched control
through much of the early and middle network, but this relationship reverses at
later depth. Across all three checkpoints, purpose disruption under the
structured intervention peaks around mid-depth and then declines toward the
output, while the matched-control trajectory remains higher after the crossover.
The precise control trajectory differs across models---including a late plateau
in Llama-3.1-70B---but the late reversal is shared across checkpoints.

These results distinguish the matched-format-associated subspace functionally
from an arbitrary perturbation of the same magnitude. They do not show that the
subspace determines generated behavior, nor that it represents format alone.
Rather, they show that intervening on this component during the actual forward
pass causally alters a downstream internal format readout, with relative
preservation of the purpose readout emerging only at later depth.

\begin{figure*}[t]
    \centering
    \includegraphics[width=\textwidth]{figures/interchange_patch_all_models.pdf}
    \caption{
        \textbf{The matched-format subspace is functionally distinguished during
        the forward computation.}
        Rank-16 interchange patching in Llama-3.1-8B swaps the matched-format
        subspace between paired benchmark and casual representations during real
        forward passes. The left panel measures whether the downstream format
        readout moves toward the matched counterpart; the right panel measures
        collateral movement of the purpose readout. Structured interchange is
        compared with random and perturbation-magnitude-matched controls.
    }
    \label{fig:interchange}
\end{figure*}


}









% ============================================================
% APPENDIX
% ============================================================
\begingroup
\color{blue}
\section{Appendix}
\label{app:methods}

\subsection{Payload generation and rendering}
\label{app:data}

We generated 320 purpose-neutral multiple-choice payloads across
16 topics. For each topic, Grok-4.3 and Gemini-3.1-Pro each generated
10 payloads. Every payload contained a context, a question, one correct
answer, and three plausible but incorrect alternatives. GPT-5.4
independently checked answer uniqueness, distractor validity, and
whether the payload could be rendered consistently across formats.
Payloads that failed validation were excluded or regenerated.

Each payload was placed into 12 template slots: six benchmark-style
templates, five casual templates, and one legacy template. The legacy
slot alternated between the original benchmark and casual templates
across payloads. Each slot produced three renderings: one stating an
evaluation purpose, one stating a deployment purpose, and one containing
no purpose statement. This produced
\[
320 \times 12 \times 3 = 11{,}520
\]
rows. The primary analysis used the 7,680 stated-purpose rows, which
form 3,840 matched evaluation/deployment pairs. The remaining 3,840
no-cue rows were analyzed separately as a control.

Response contracts and answer positions were assigned independently of
purpose. This prevents the evaluation and deployment labels from being
systematically associated with a particular answer location or requested
response style.

\subsection{Purpose-vocabulary families}
\label{app:vocabulary}

Purpose statements were divided into five families with disjoint
purpose-specific vocabulary. The families expressed the same
evaluation-versus-deployment distinction through different settings:

\begin{table}[t]
\centering
\small
\begin{tabular}{p{0.09\linewidth}p{0.38\linewidth}p{0.38\linewidth}}
\toprule
Family & Evaluation setting & Deployment setting \\
\midrule
1 & Editorial review, using terms such as reviewer, review board,
critique, and go over
& Purchasing, using terms such as customer, purchasing crew,
act on, and draw on \\

2 & Exam grading, using terms such as grader, examiners,
mark, and grade
& Client services, using terms such as client, callers,
follow, and heed \\

3 & Audit or inspection, using terms such as auditor, audit unit,
inspect, and comb through
& Tenant maintenance, using terms such as tenant, maintenance staff,
apply, and carry out \\

4 & Scored testing, using terms such as scorer, testing lab,
rate, and score
& User support, using terms such as user, support line,
employ, and forward \\

5 & Judged verification, using terms such as judge, verifiers,
appraise, and vet
& Household shopping, using terms such as shopper, householders,
adopt, and put into practice \\
\bottomrule
\end{tabular}
\caption{Settings and representative vocabulary used by the five
purpose-statement families.}
\label{tab:vocab-families}
\end{table}

Each family contained six matched evaluation/deployment statement
pairs. Two pairs varied both the relevant nouns and verbs, two varied
only the nouns, and two varied only the verbs. The main dataset used
the stated third-person versions of these statements.

Each payload was assigned three of the five families using a
deterministic, topic-balanced schedule. Across its 12 template slots,
each assigned family appeared four times. Consequently, every family
contained 1,536 stated-purpose rows, and every topic was represented
in every held-out family.

\subsection{Activation extraction}
\label{app:activations}

We extracted activations from the pinned instruction-tuned checkpoints
of Llama-3.1-8B, Llama-3.1-70B, and Llama-3.3-70B. The corresponding
models contain 32, 80, and 80 transformer blocks.

Each rendered prompt was passed to the checkpoint's chat template as a
single user message with the generation prompt enabled. At every
zero-indexed transformer block, we recorded the hidden state at the
final prompt token. Activations were collected after the block output
and before the model's final normalization layer. Inference used BF16,
and the extracted activations were stored in FP32. Checkpoint revisions,
dataset hashes, activation hashes, and extraction settings are provided
in the reproducibility inventory.

\subsection{Probe fitting and evaluation}
\label{app:probes}

At each layer, a probe consisted of a standardizer followed by
$L_2$-regularized logistic regression. We used $C=0.1$, a maximum of
3,000 iterations, and random state 0. Both the standardizer and the
classifier were fitted only on the training rows in each fold.

We evaluated probes using five leave-one-family-out folds. In each
fold, four purpose-vocabulary families were used for training and the
remaining family was used for testing. The same payload could occur in
training and testing under different families. Thus, this split tests
generalization to unseen purpose wording, not generalization to unseen
payloads.

Every stated-purpose row received one out-of-fold prediction. We then
combined these predictions across the five folds and calculated one ROC
AUC and one balanced-accuracy value for each layer and probe regime.
Balanced accuracy used the fitted classifier's zero decision threshold.

The probe regimes were defined as follows:

\begin{table}[t]
\centering
\small
\begin{tabular}{p{0.23\linewidth}p{0.29\linewidth}p{0.34\linewidth}}
\toprule
Probe & Training rows & Evaluation rows \\
\midrule
Confounded
& Evaluation in benchmark format and deployment in casual format
& All four purpose-by-format cells \\

Decorrelated
& All four purpose-by-format cells
& All four purpose-by-format cells \\

Benchmark-to-casual
& Benchmark-formatted rows
& Casual-formatted rows \\

Casual-to-benchmark
& Casual-formatted rows
& Benchmark-formatted rows \\

Format
& All stated-purpose rows, labeled by format
& All stated-purpose rows \\
\bottomrule
\end{tabular}
\caption{Training and evaluation data used by each probe regime.}
\label{tab:probe-regimes}
\end{table}

\paragraph{Surface-text control.}
We removed each explicit purpose statement using its recorded character
span. After removal, the evaluation and deployment versions of a
matched pair were byte-identical. We then fitted logistic-regression
classifiers using word and character TF--IDF features, holding out
templates during evaluation. A stricter version held out both the
template and the purpose-vocabulary family. This analysis checks for
purpose-label leakage elsewhere in the rendered prompt. It does not
determine whether the activation-level signal depends on the lexical
or semantic content of the purpose statement itself.

\paragraph{Equal-sample-size control.}
The decorrelated probe normally receives more training rows than either
single-format probe. We therefore repeated the comparison at
Llama-3.1-8B layer 12 while holding the number of training rows constant.
For each held-out family, the benchmark-to-casual and
casual-to-benchmark probes selected one format from every training
payload block. The decorrelated condition also selected one format per
block, with half of the blocks assigned to each format. We repeated this
assignment using four fixed random seeds.

\subsection{Bootstrap uncertainty}
\label{app:bootstrap}

We calculated uncertainty intervals by resampling complete payload
blocks rather than individual rows. All purpose, format, template, and
cue-family renderings belonging to a sampled payload remained together.
This preserves the matched factorial structure and avoids treating
renderings of the same underlying task as independent observations.

When comparing layers or probe conditions, the same bootstrap draw was
used across the quantities being compared. Layers were treated as
repeated observations along a model's depth trajectory, not as
independent experimental replicates.

\subsection{Format-subspace estimation and projection}
\label{app:subspace}

For each complete matched group, we calculated the format difference
\[
B =
\frac{
(h_{EB}+h_{DB})-(h_{EC}+h_{DC})
}{4}.
\]
At each layer, the resulting $B$ vectors were placed into a matrix and
decomposed using singular value decomposition. The leading right
singular vectors defined the estimated format subspace. We examined
subspace ranks 1, 4, 16, and 64.

Subspace energy was compared with the isotropic expectation $r/d$,
where $r$ is the subspace rank and $d$ is the activation dimension.
We also generated 10,000 rank-matched random subspaces within the
leading 64-dimensional empirical activation subspace.

Projection analyses were cross-fitted. For each vocabulary-family fold,
the format basis was estimated using only the training groups.
Training and held-out activations were then projected as
\[
h'=(I-QQ^\top)h,
\]
where the columns of $Q$ form the estimated format basis. The
standardizer and probe were refitted using the projected training
activations before evaluation on the projected held-out activations.

\subsection{Interchange and prospective interventions}
\label{app:interventions}

For the interchange analysis, we replaced the format-subspace component
of an activation with the corresponding component from its matched
opposite-format prompt:
\[
h_i' =
h_i - QQ^\top h_i + QQ^\top h_j,
\]
where $i$ and $j$ share the same payload and stated purpose but use
different response formats. The modified state was inserted at the
final prompt token after the selected transformer block, and the
remaining blocks were allowed to run normally.

The primary interchange analysis used a rank-16 format subspace at
every layer of Llama-3.1-8B. Rank-matched random-subspace swaps and
equal-norm random perturbations were used as controls. Readouts and
subspaces were fitted only from the training portion of each
vocabulary-family fold. Confidence intervals and the joint
worst-layer statistic used complete-payload bootstrap resampling.

For the prospective interaction intervention, we first removed from
the estimated interaction direction $C$ any component lying in the
span of the purpose and format main effects:
\[
C_{\perp} =
\left(I-P_{[A,B]}\right)C.
\]
Directions, intervention scales, readouts, source and target layers,
transfer payloads, and random controls were fixed before outcome
evaluation. We applied sign-symmetric positive and negative
interventions and measured the high-minus-low change in the frozen
interaction readout. Half-dose interventions and matched random
directions served as additional controls. Statistical testing used
payload-block bootstrap intervals, Holm correction across the two
preregistered models, and a frozen random-control gate. Results from
the third checkpoint were treated as an extension and were not pooled
with the preregistered tests.

\endgroup